% Common preamble for all lab reports
\documentclass[a4paper, 12pt]{article}

% Language and encoding
\usepackage[T2A]{fontenc}
\usepackage[utf8]{inputenc}
\usepackage[ukrainian]{babel}

% Typography
\usepackage{microtype}
\usepackage{lmodern}

% Page geometry
\usepackage[margin=25mm, headheight=14pt]{geometry}

% Math
\usepackage{amsmath, amssymb}

% Graphics and figures
\usepackage{graphicx}
\usepackage{float}
\usepackage{subcaption}
\usepackage[justification=centering]{caption}

% Tables
\usepackage{tabularx}
\usepackage{booktabs}
\usepackage{array}

% Colors and hyperlinks
\usepackage{xcolor}
\usepackage[colorlinks=true, linkcolor=blue!60!black, urlcolor=blue!60!black, citecolor=blue!60!black]{hyperref}

% Lists
\usepackage{enumitem}
\setlist{nosep, leftmargin=*}

% Section formatting
\usepackage{titlesec}
\titleformat{\section}{\large\bfseries}{\thesection.}{0.5em}{}
\titleformat{\subsection}{\normalsize\bfseries}{\thesubsection.}{0.5em}{}

% Header/footer
\usepackage{fancyhdr}
\pagestyle{fancy}
\fancyhf{}
\fancyhead[L]{\small\leftmark}
\fancyhead[R]{\small\thepage}
\renewcommand{\headrulewidth}{0.4pt}
\fancypagestyle{plain}{\fancyhf{}\fancyfoot[C]{\thepage}\renewcommand{\headrulewidth}{0pt}}

% Placeholder image command
\newcommand{\placeholder}[2][0.8\textwidth]{%
  \begin{figure}[H]
    \centering
    \fbox{\parbox{#1}{\vspace{4cm}\centering\textit{#2}\vspace{4cm}}}
    \caption{#2}
  \end{figure}
}

% Placeholder for inline result values
\newcommand{\result}[1]{\textbf{\textcolor{red!70!black}{[#1]}}}

\usepackage{../common/labstyle}

\labnum{3}
\labtitle{Теплообмiн}

\begin{document}
\maketitle
\tableofcontents
\newpage

% ===================================================================
\section{Вихiднi данi}
% ===================================================================

\begin{table}[H]
  \centering
  \caption{Вихiднi данi (варiант~13)}
  \begin{tabularx}{0.80\textwidth}{Xr}
    \toprule
    \textbf{Параметр} & \textbf{Значення} \\
    \midrule
    Верхня кришка                              & Плоска \\
    Нижня кришка                               & Сферична \\
    Товщина стiнки ємностi                     & 2{,}1~мм \\
    Температура гарячого контуру               & 108\textdegree C (381{,}15~K) \\
    Температура холодної рiдини                & 12\textdegree C (285{,}15~K) \\
    Швидкiсть потоку в контурi                 & 1{,}4~м/с \\
    Швидкiсть потоку в ємностi                 & 0{,}006~м/с \\
    Початковий тиск                            & 1~атм (101325~Па) \\
    Матерiал контуру (тверде тiло)             & Copper \\
    Рiдина                                     & Water \\
    Модель турбулентностi                      & $k$-$\varepsilon$ \\
    \bottomrule
  \end{tabularx}
\end{table}

Геометрiя ємностi --- з Практичної роботи~\textnumero 1: цилiндр ($H=1000$~мм, $D=500$~мм), плоска верхня кришка, сферична нижня кришка.

\subsection{Ключовi формули}

Закон Фур'є: $q = -\lambda\,\partial T / \partial x$.
Рiвняння енергiї: $\rho c_p(\partial T/\partial t + \mathbf{U}\cdot\nabla T) = \nabla\cdot(\lambda\nabla T) + S_E$.
Умова спряження (CHT) на iнтерфейсi Fluid--Solid:
\begin{equation}
  T_{\mathrm{solid}}\big|_{\mathrm{int}} = T_{\mathrm{fluid}}\big|_{\mathrm{int}},
  \qquad
  q_{\mathrm{solid}}\big|_{\mathrm{int}} = q_{\mathrm{fluid}}\big|_{\mathrm{int}}.
\end{equation}

Тепловий потiк через контур:
\begin{equation}
  \dot{Q} = \dot{m}\,c_p\,\Delta T.
\end{equation}

% ===================================================================
\section{Послiдовнiсть виконання}
% ===================================================================

Робота виконана через ANSYS Workbench з солвером CFX (Conjugate Heat Transfer).

\begin{lstlisting}[style=scheme]
1. Geometry: CadQuery (geometry/lab3.py + common.py)
   |-- make_vessel_inner_volume(flat_bottom=True):
   |   cylinder D_in=495.8mm, H=1000mm
   |-- Coil: 6 vertical pipes at R=130mm, 60 deg apart,
   |   OD=25mm, wall=2mm, connected by U-bends
   |-- Boolean: vessel - outer_coil -> vessel_fluid
   |-- Assembly (3 bodies): vessel_fluid,
   |   coil_solid_copper, coil_fluid
   |-- Export STEP -> geometry/step/lab3.step
2. ANSYS Workbench: CreateSystem("Fluid Flow (CFX)")
3. Import lab3.step -> SpaceClaim (3 bodies)
4. Enter Meshing (Python via SendCommand)
   |-- ID bodies by volume: vessel > solid > fluid
   |-- NS "volume_inlet"  -> vessel bottom face (Z=0)
   |-- NS "volume_outlet" -> vessel top face (Z=1000)
   |-- NS "Vessel_Wall"   -> remaining vessel faces
   |-- NS "coil_inside_inlet"  -> coil pipe 0 at Z=1000
   |-- NS "coil_inside_outlet" -> coil pipe 5 at Z=1000
   |-- Mesh: 5mm -> 384709 nodes, 1990506 elements
5. CFX-Pre (CCL import)
   |-- Domain "volume": Water, Total Energy, k-eps
   |     Inlet: 285.15 K, 0.006 m/s
   |     Outlet: Avg Static Pressure 0 Pa
   |-- Domain "pipe" (solid): Copper, Thermal Energy
   |-- Domain "pipe_inside": Water, Total Energy, k-eps
   |     Inlet: 381.15 K, 1.4 m/s
   |     Outlet: Avg Static Pressure 0 Pa
   |-- Interfaces: GGI Fluid-Solid CHT (auto-connect)
   |-- Solver: High Resolution, 100 iter, RMS 1e-4
6. Solve -> 100 iterations (12 cores)
7. CFX-Post (Power Syntax)
   |-- Cross-section planes (XZ, XY)
   |-- Temperature contours (vessel + coil domains)
   |-- Velocity streamlines (both fluid domains)
   |-- Function Calculator: ave(T), ave(V), ave(P)
   |   on all 4 inlets/outlets
   |-- Export images (PNG) + CSV
\end{lstlisting}

\subsection{Геометрiя (CadQuery $\to$ STEP)}

Скрипт \texttt{geometry/lab3.py} будує збiрку з трьох тiл:
\begin{itemize}
  \item \textbf{vessel\_fluid} --- внутрiшнiй об'єм ємностi (цилiндр $D_{\mathrm{in}}=495{,}8$~мм, $H=1000$~мм з плоским днищем);
  \item \textbf{coil\_solid\_copper} --- тiло мiдного змiйовика (труба $D_{\mathrm{out}}=25$~мм, стiнка~2~мм): 6 вертикальних труб на радiусi $R=130$~мм з напiвкруглими U-подiбними зворотами;
  \item \textbf{coil\_fluid} --- внутрiшнiй об'єм змiйовика ($D_{\mathrm{in}}=21$~мм).
\end{itemize}

\begin{figure}[H]
  \centering
  \includegraphics[width=0.65\textwidth]{img/geometry}
  \caption{Геометрiя ємностi з iнтегрованим трубчастим контуром (6 вертикальних труб з U-зворотами)}
\end{figure}

Обидва кiнцi змiйовика виходять через верхню кришку ємностi. Рiдинний об'єм ємностi отримано як Boolean Subtract вихiдного цилiндра мiнус зовнiшня поверхня труб.

% ===================================================================
\section{Вiзуалiзацiя розрахункової моделi}
% ===================================================================

\begin{figure}[H]
  \centering
  \includegraphics[width=0.65\textwidth]{img/mesh_overview}
  \caption{Розрахункова сiтка (Mesh), загальний вигляд ємностi}
\end{figure}

\begin{figure}[H]
  \centering
  \includegraphics[width=0.65\textwidth]{img/mesh_closeup}
  \caption{Ущiльнення сiтки навколо труб контуру (поперечний перерiз)}
\end{figure}

\begin{figure}[H]
  \centering
  \includegraphics[width=0.55\textwidth]{img/mesh_pipe_detail}
  \caption{Детальний вигляд сiтки навколо однiєї труби контуру (тверде тiло --- оранжевий, рiдина --- блакитний)}
\end{figure}

Сiтка побудована з глобальним розмiром елемента 5~мм. Загальна кiлькiсть вузлiв --- 384709, елементiв --- 1990506 (тетраедри) у трьох доменах:
\begin{itemize}
  \item \texttt{volume} (рiдина ємностi): 302238 вузлiв, 1701959 елементiв;
  \item \texttt{pipe} (мiдна стiнка контуру): 44444 вузли, 136086 елементiв;
  \item \texttt{pipe\_inside} (рiдина контуру): 38027 вузлiв, 152461 елемент.
\end{itemize}

На детальному виглядi видно ущiльнення сiтки у зонi контакту Fluid--Solid: елементи мiдної стiнки (оранжевий) та рiдини ємностi (блакитний) узгодженi через GGI-iнтерфейс.

% ===================================================================
\section{Результати розрахунку}
% ===================================================================

\subsection{Temperature}

\begin{figure}[H]
  \centering
  \includegraphics[width=0.85\textwidth]{img/temp_cross_section}
  \caption{Розподiл температури на поперечному перерiзi --- домен ємностi (VolumeC) та контуру (CoilC), дiапазон 285{,}15--290{,}15~K}
\end{figure}

На поперечному перерiзi видно локальнi зони нагрiву навколо кожної з 6 труб контуру. Температура рiдини в ємностi пiдвищується вiд 285{,}15~K (12\textdegree C) поблизу стiнок до 290{,}15~K (17\textdegree C) безпосередньо навколо труб. Нагрiв концентрується у безпосереднiй близькостi до поверхнi контуру, що свiдчить про переважно кондуктивний механiзм теплопередачi при низькiй швидкостi потоку в ємностi.

\begin{figure}[H]
  \centering
  \includegraphics[width=0.65\textwidth]{img/temp_coil}
  \caption{Розподiл температури на поздовжньому перерiзi контуру (дiапазон 285{,}15--297{,}7~K)}
\end{figure}

На поздовжньому перерiзi контуру видно поступове охолодження гарячої рiдини вздовж труб: максимальна температура (297{,}7~K) на входi, мiнiмальна (285{,}15~K) --- у зонах, де контур контактує з холодною рiдиною ємностi. Градiєнт температури бiльш виражений у зонах U-зворотiв.

\subsubsection{Чисельнi значення температури}

\begin{table}[H]
  \centering
  \caption{Середня температура на входах та виходах доменiв}
  \begin{tabularx}{0.85\textwidth}{Xrrr}
    \toprule
    \textbf{Границя} & \textbf{$T$, K} & \textbf{$T$, \textdegree C} & \textbf{$\Delta T$, K} \\
    \midrule
    Контур --- вхiд (coil\_inside\_inlet)   & 381{,}150  & 108{,}00  & --- \\
    Контур --- вихiд (coil\_inside\_outlet) & 380{,}699  & 107{,}55  & 0{,}451 \\
    \midrule
    Ємнiсть --- вхiд (volume\_inlet)        & 285{,}150  & 12{,}00   & --- \\
    Ємнiсть --- вихiд (volume\_outlet)      & 285{,}617  & 12{,}47   & 0{,}467 \\
    \bottomrule
  \end{tabularx}
\end{table}

Гаряча рiдина в контурi охолоджується на $\Delta T_{\mathrm{coil}} = 0{,}451$~K, холодна рiдина в ємностi нагрiвається на $\Delta T_{\mathrm{vessel}} = 0{,}467$~K.

\subsection{Velocity}

\begin{figure}[H]
  \centering
  \includegraphics[width=0.95\textwidth]{img/velocity_streamlines}
  \caption{Лiнiї току (Streamlines), зафарбованi за швидкiстю --- домен ємностi (VolumeS) та контуру (CoilS), дiапазон 0--1{,}568~м/с}
\end{figure}

На зображеннi видно два характернi режими потоку:
\begin{itemize}
  \item \textbf{Контур (CoilS)}: швидкiсть ~1{,}4~м/с, потiк рухається через вертикальнi труби з прискоренням у U-зворотах (до 1{,}568~м/с). Лiнiї току щiльнi та упорядкованi --- турбулентний потiк у трубi.
  \item \textbf{Ємнiсть (VolumeS)}: швидкiсть ~0{,}006~м/с, рiвномiрний ламiнарний потiк знизу вгору. Лiнiї току паралельнi, обтiкають труби контуру.
\end{itemize}

\subsubsection{Чисельнi значення швидкостi}

\begin{table}[H]
  \centering
  \caption{Середня швидкiсть на входах та виходах доменiв}
  \begin{tabularx}{0.80\textwidth}{Xr}
    \toprule
    \textbf{Границя} & \textbf{$V$, м/с} \\
    \midrule
    Контур --- вхiд (coil\_inside\_inlet)   & 1{,}397 \\
    Контур --- вихiд (coil\_inside\_outlet) & 1{,}374 \\
    \midrule
    Ємнiсть --- вхiд (volume\_inlet)        & 0{,}00600 \\
    Ємнiсть --- вихiд (volume\_outlet)      & 0{,}00592 \\
    \bottomrule
  \end{tabularx}
\end{table}

Швидкiсть на виходi контуру (1{,}374~м/с) дещо менша за вхiдну (1{,}397~м/с) через гiдравлiчнi втрати у U-зворотах та тертi об стiнки труб.

\subsection{Тиск на входах та виходах}

\begin{table}[H]
  \centering
  \caption{Середнiй тиск на входах та виходах доменiв}
  \begin{tabularx}{0.80\textwidth}{Xrr}
    \toprule
    \textbf{Границя} & \textbf{$p$, Па} & \textbf{$\Delta p$, Па} \\
    \midrule
    Контур --- вхiд   & 111229 & --- \\
    Контур --- вихiд  & 101335 & 9894 \\
    \midrule
    Ємнiсть --- вхiд  & 101325 & --- \\
    Ємнiсть --- вихiд & 101325 & $\approx 0$ \\
    \bottomrule
  \end{tabularx}
\end{table}

Перепад тиску в контурi $\Delta p_{\mathrm{coil}} = 9894$~Па ($\approx 9{,}9$~кПа) обумовлений значною довжиною трубопроводу (6 вертикальних труб + 5 U-зворотiв) та високою швидкiстю потоку. В ємностi перепад тиску практично вiдсутнiй через низьку швидкiсть.

% ===================================================================
\section{Аналiз та висновок}
% ===================================================================

\subsection{Тепловий баланс}

Оцiнка теплового потоку через контур:
\begin{itemize}
  \item Масова витрата контуру: $\dot{m}_{\mathrm{coil}} = \rho \cdot V \cdot A = 997 \cdot 1{,}4 \cdot \pi \cdot 0{,}0105^2 \approx 0{,}483$~кг/с.
  \item Тепловий потiк контуру: $\dot{Q}_{\mathrm{coil}} = \dot{m}_{\mathrm{coil}} \cdot c_p \cdot \Delta T_{\mathrm{coil}} = 0{,}483 \cdot 4182 \cdot 0{,}451 \approx 910$~Вт.
\end{itemize}

Оцiнка теплового потоку через ємнiсть:
\begin{itemize}
  \item Масова витрата ємностi: $\dot{m}_{\mathrm{vessel}} = 997 \cdot 0{,}006 \cdot \pi \cdot 0{,}2479^2 \approx 1{,}156$~кг/с.
  \item Тепловий потiк ємностi: $\dot{Q}_{\mathrm{vessel}} = 1{,}156 \cdot 4182 \cdot 0{,}467 \approx 2258$~Вт.
\end{itemize}

З даних солвера, тепловий потiк через стiнку контуру складає $\dot{Q}_{\mathrm{wall}} \approx 229$~Вт (pipe\_inside $\to$ pipe: 229{,}2~Вт, pipe $\to$ volume: 228{,}3~Вт). Розбiжнiсть з оцiнкою за $\Delta T$ пояснюється тим, що розрахунок обмежений 100 iтерацiями --- тепловi рiвняння ще не досягли повної збiжностi (залишки рiвняння енергiї на рiвнi $10^{-4}$--$10^{-3}$, iмбаланс ~0{,}1\%).

Масовi витрати: контур $\dot{m}_{\mathrm{coil}} = 0{,}477$~кг/с, ємнiсть $\dot{m}_{\mathrm{vessel}} = 6{,}92 \times 10^{-3}$~кг/с. Iмбаланс мас менше 0{,}001\% для обох доменiв.

\subsection{Ефективнiсть теплообмiну}

Ефективнiсть теплообмiнника:
\[
  \varepsilon = \frac{\dot{Q}_{\mathrm{actual}}}{\dot{Q}_{\mathrm{max}}}
  = \frac{\dot{Q}_{\mathrm{wall}}}{\dot{m}_{\min} \cdot c_p \cdot (T_{\mathrm{hot,in}} - T_{\mathrm{cold,in}})}
  = \frac{229}{6{,}92 \times 10^{-3} \cdot 4182 \cdot 96}
  = \frac{229}{2779}
  \approx 0{,}082 \; (8{,}2\%).
\]

Помiрна ефективнiсть ($\varepsilon \approx 8{,}2\%$) обумовлена:
\begin{itemize}
  \item Високою швидкiстю гарячої рiдини в контурi (1{,}4~м/с) --- малий час перебування рiдини в зонi теплообмiну.
  \item Обмеженою поверхнею контакту --- 6 прямих труб замiсть спiрального змiйовика.
  \item Недостатньою збiжнiстю розрахунку (100 iтерацiй).
\end{itemize}

\subsection{Фiзична реалiстичнiсть}

Результати моделювання демонструють фiзично реалiстичну картину:
\begin{itemize}
  \item Гаряча рiдина вiддає тепло через мiдну стiнку контуру холоднiй рiдинi в ємностi --- напрямок теплопередачi правильний.
  \item Нагрiв холодної рiдини локалiзований навколо труб контуру, що вiдповiдає кондуктивно-конвективному механiзму при низькiй швидкостi потоку.
  \item Перепад тиску в контурi ($\approx 9{,}9$~кПа) фiзично обґрунтований для труби дiаметром 21~мм при швидкостi 1{,}4~м/с та загальнiй довжинi ~5~м.
  \item Температурнi градiєнти не перевищують розумних меж --- немає нефiзичних стрибкiв.
  \item Умова CHT (спряження) коректно забезпечує безперервнiсть температури та теплового потоку на iнтерфейсах Fluid--Solid.
\end{itemize}

\textbf{Висновок:} Моделювання сполученого теплообмiну (CHT) у системi <<ємнiсть~--- трубчастий контур>> при заданих умовах (варiант~13) демонструє фiзично коректну картину: гаряча рiдина ($T=108$\textdegree C, $V=1{,}4$~м/с) в мiдному контурi вiддає тепло холоднiй рiдинi ($T=12$\textdegree C, $V=0{,}006$~м/с) в ємностi. Температура на виходi контуру знизилась до 107{,}55\textdegree C ($\Delta T = 0{,}45$~K), температура на виходi ємностi зросла до 12{,}47\textdegree C ($\Delta T = 0{,}47$~K). Тепловий потiк через стiнку контуру $\dot{Q} \approx 229$~Вт. Для пiдвищення ефективностi рекомендується збiльшити кiлькiсть iтерацiй, зменшити швидкiсть потоку в контурi або збiльшити площу теплообмiнної поверхнi (наприклад, використати спiральний змiйовик замiсть прямих труб).

% ===================================================================
\appendix
\section{Скрипт генерацiї геометрiї}
% ===================================================================

\subsection{common.py --- спiльна геометрiя ємностi}

\begin{lstlisting}[style=python]
"""Shared vessel geometry for Labs 1, 3, 5 (Variant 13)."""

import cadquery as cq

OD = 500.0
WALL = 2.1
R_IN = (OD - 2 * WALL) / 2
SHELL_H = 1000.0


def make_vessel_inner_volume(flat_bottom=False):
    inner_cyl = cq.Workplane("XY").circle(R_IN).extrude(SHELL_H)
    if flat_bottom:
        return inner_cyl

    inner_sphere = cq.Workplane("XY").sphere(R_IN)
    cut_box = (cq.Workplane("XY")
        .box(OD+10, OD+10, OD+10, centered=True)
        .translate((0, 0, (OD+10)/2)))
    bottom_hemi = inner_sphere.cut(cut_box)
    return inner_cyl.union(bottom_hemi)
\end{lstlisting}

\subsection{lab3.py --- трубчастий теплообмiнник}

\begin{lstlisting}[style=python]
"""Lab 3 -- CFX CHT geometry, Variant 13.

Vertical pipe-bundle heat exchanger inside a vessel:
  1. vessel_fluid -- cylinder, flow bottom->top
  2. coil_solid_copper -- copper pipe wall (6 pipes + U-bends)
  3. coil_fluid -- fluid inside the pipe
"""

import math
from pathlib import Path
import cadquery as cq
from cadquery import Vector
from OCP.BRepOffsetAPI import BRepOffsetAPI_MakePipeShell
from common import make_vessel_inner_volume, R_IN, SHELL_H

TUBE_OD = 25.0
TUBE_WALL = 2.0
TUBE_OR = TUBE_OD / 2.0
TUBE_IR = (TUBE_OD - 2 * TUBE_WALL) / 2.0

N_PIPES = 6
R_COIL = 130.0
Z_BOTTOM = 80.0
Z_TOP = 880.0
Z_EXIT = SHELL_H
PIPE_ANGLES = [math.radians(i * 360 / N_PIPES)
               for i in range(N_PIPES)]


def _pipe_pos(i):
    a = PIPE_ANGLES[i]
    return R_COIL * math.cos(a), R_COIL * math.sin(a)


def _make_uturn(x1, y1, x2, y2, z_level, goes_below):
    mx, my = (x1+x2)/2, (y1+y2)/2
    r = math.sqrt((x2-x1)**2 + (y2-y1)**2) / 2
    if goes_below:
        mid = Vector(mx, my, z_level - r)
    else:
        mid = Vector(mx, my, z_level + r)
    return cq.Edge.makeThreePointArc(
        Vector(x1, y1, z_level), mid,
        Vector(x2, y2, z_level))


def _make_coil_wire():
    edges = []
    x0, y0 = _pipe_pos(0)
    edges.append(cq.Edge.makeLine(
        Vector(x0, y0, Z_EXIT),
        Vector(x0, y0, Z_BOTTOM)))
    for i in range(N_PIPES - 1):
        x1, y1 = _pipe_pos(i)
        x2, y2 = _pipe_pos(i + 1)
        going_down = (i % 2 == 0)
        if going_down:
            edges.append(_make_uturn(
                x1, y1, x2, y2, Z_BOTTOM, True))
            edges.append(cq.Edge.makeLine(
                Vector(x2, y2, Z_BOTTOM),
                Vector(x2, y2, Z_TOP)))
        else:
            edges.append(_make_uturn(
                x1, y1, x2, y2, Z_TOP, False))
            edges.append(cq.Edge.makeLine(
                Vector(x2, y2, Z_TOP),
                Vector(x2, y2, Z_BOTTOM)))
    xl, yl = _pipe_pos(N_PIPES - 1)
    last_down = ((N_PIPES - 2) % 2 != 0)
    z_last = Z_BOTTOM if last_down else Z_TOP
    edges.append(cq.Edge.makeLine(
        Vector(xl, yl, z_last),
        Vector(xl, yl, Z_EXIT)))
    return cq.Wire.assembleEdges(edges)


def _sweep(radius, wire):
    start = wire.startPoint()
    tangent = wire.tangentAt(0)
    circle = cq.Wire.makeCircle(
        radius, center=start, normal=tangent)
    builder = BRepOffsetAPI_MakePipeShell(wire.wrapped)
    builder.SetMode(True)  # Frenet
    builder.Add(circle.wrapped)
    builder.Build()
    builder.MakeSolid()
    solid = cq.Solid(builder.Shape())
    if solid.Volume() < 0:
        solid = cq.Solid(solid.wrapped.Reversed())
    return solid


def build():
    wire = _make_coil_wire()
    outer = _sweep(TUBE_OR, wire)
    inner = _sweep(TUBE_IR, wire)
    outer_wp = cq.Workplane("XY").newObject([outer])
    inner_wp = cq.Workplane("XY").newObject([inner])
    coil_solid = outer_wp.cut(inner_wp)
    coil_fluid = inner_wp
    vessel_fluid = make_vessel_inner_volume(flat_bottom=True)
    vessel_fluid = vessel_fluid.cut(outer_wp)
    return vessel_fluid, coil_solid, coil_fluid


vessel_fluid, coil_solid, coil_fluid = build()

assy = cq.Assembly()
assy.add(vessel_fluid, name="vessel_fluid")
assy.add(coil_solid, name="coil_solid_copper")
assy.add(coil_fluid, name="coil_fluid")

out_dir = Path(__file__).parent / "step"
out_dir.mkdir(exist_ok=True)
assy.export(str(out_dir / "lab3.step"))
\end{lstlisting}

\end{document}
